%% eeprep.com resume template (single column, one page, pdflatex)
%% Fictional student: Alex Rivera. NOT A REAL PERSON. Every line matches
%% content-plans/resume.md. Replace it all with your own.
%% Based on Edward Silva's resume advice and the r/EngineeringResumes wiki:
%% https://www.reddit.com/r/EngineeringResumes/wiki/index
\documentclass[letterpaper,11pt]{article}

%% Rule: machine-readable text. glyphtounicode maps every glyph to a Unicode
%% character, so applicant tracking systems and AI readers parse clean text.
\input{glyphtounicode}
\pdfgentounicode=1

\usepackage[T1]{fontenc}   % real glyphs for accents and dashes, no odd ligature splits
\usepackage{lmodern}       % scalable Latin Modern fonts
\usepackage[margin=0.5in]{geometry}
\usepackage{enumitem}
\usepackage{titlesec}
\usepackage[hidelinks]{hyperref}

%% Rule: PDF metadata. Anyone can see it under document properties. It costs
%% nothing, but do not count on hiring software reading it.
\hypersetup{
  pdftitle    = {Alex Rivera - Resume},
  pdfauthor   = {Alex Rivera},
  pdfsubject  = {Electrical engineering internship resume},
  pdfkeywords = {electrical engineering, power electronics, PCB design, KiCad, LTspice, STM32, embedded C, FPGA, Verilog, Python}
}

%% Rule: consistent spacing. Every gap between sections comes from this one
%% length. Change it here, never with a one-off \vspace in the body.
\newlength{\secgap}
\setlength{\secgap}{6pt}
\newlength{\itemgap}           % gap between entries inside a section
\setlength{\itemgap}{4pt}

\pagestyle{empty}              % no page number on a one-page resume
\setlength{\parindent}{0pt}    % Rule: left-justify everything
\raggedright
\urlstyle{same}

%% Section headings: bold, a thin rule, no centering.
\titleformat{\section}{\large\bfseries}{}{0pt}{}[\vspace{-6pt}\rule{\textwidth}{0.4pt}]
\titlespacing*{\section}{0pt}{\secgap}{\itemgap}

%% Tight bullets that use the full line width.
\setlist[itemize]{leftmargin=1.2em, topsep=1pt, itemsep=0pt, parsep=0pt, partopsep=0pt}

%% \entry{bold title}{right-aligned dates}
%% Rule: dates go flush right with \hfill so the eye scans titles on the left.
\newcommand{\entry}[2]{\textbf{#1}\hfill #2\par}
%% \skillsused{...}
%% Rule: put the skills you used under the project title. Repeating them in
%% the Skills section is good, because it shows where you used each one.
\newcommand{\skillsused}[1]{\textit{#1}\par}
%% Header separator. The contact line is set at 9.5 pt so all six items fit
%% on one line, as the persona and Edward's advice ask.
\newcommand{\sep}{\hspace{0.35em}|\hspace{0.35em}}

\begin{document}

%% ---------------- Header ----------------
%% Rules: name on top, no line under it, no logos, no "Phone:" or "Email:"
%% labels. Write the full LinkedIn URL as text. Everything fits on one line
%% under the name, including citizenship and relocation.
{\LARGE\bfseries Alex Rivera}\par
\vspace{2pt}
{\fontsize{9.5}{11.5}\selectfont (555) 555-0142\sep arivera@example.com\sep
\href{https://linkedin.com/in/alex-rivera-example}{linkedin.com/in/alex-rivera-example}\sep
Columbus, OH\sep US citizen\sep Willing to relocate\par}

%% Rule: no objective or summary. It takes space and tells a recruiter little.
%% Rule: no interests section unless you can name specific things.

%% ---------------- Projects ----------------
%% Rule: projects come first. They are your strongest proof for an internship.
%% Rule: XYZ bullets. "Accomplished [X] as measured by [Y], by doing [Z]."
%% Every bullet has a number. Fill the line or keep it to one line.
\section{Projects}

\entry{Synchronous Buck Converter for a CubeSat Power Board}{Fall 25 - Spring 26}
\skillsused{KiCad, LTspice, STM32, Embedded C, bench testing with an electronic load}
\begin{itemize}
  \item Reached 92\% peak efficiency at a 2 A load, measured on an electronic load, by replacing the catch diode with a synchronous MOSFET stage and tuning dead time to 40 ns with an STM32 timer, up from 84\% before
  \item Cut output ripple from 85 mV to 18 mV peak to peak by redesigning the 4-layer PCB in KiCad with a solid ground plane and a tighter switching loop, then resizing the LC output filter for 500 kHz switching
  \item Wrote STM32 firmware in C that samples input current at 1 kHz and shuts the rail down within 50 \textmu s of an overcurrent fault, which kept the board alive through 3 deliberate short-circuit tests at a full 12 V input
  \item Kept the FETs under 50 \textdegree{}C at full load by adding thermal vias and a 2 oz copper pour under them
\end{itemize}
\vspace{\itemgap}

\entry{FPGA Five-Band Audio Equalizer}{Spring 25}
\skillsused{Verilog, Vivado, MATLAB, Basys 3 FPGA, fixed-point DSP}
\begin{itemize}
  \item Processed 48 kHz stereo audio at 0.4 ms latency by pipelining five 64-tap FIR filters in Verilog
  \item Matched the MATLAB response within 0.1 dB from 20 Hz to 20 kHz by testing coefficients in Vivado
  \item Cut LUT use 35\% by sharing one multiplier per filter across all 64 taps with time-division scheduling
\end{itemize}
\vspace{\itemgap}

\entry{Bluetooth Battery Monitor}{Fall 24}
\skillsused{ESP32, Arduino C++, Python, INA219 current sensor}
\begin{itemize}
  \item Tracked 4 Li-ion cells within 10 mV each with an ESP32 monitor that streams to a Python dashboard
  \item Extended pack runtime 15\% in field tests by adding a cutoff that drops non-critical loads at 3.3 V
\end{itemize}

%% ---------------- Experience ----------------
%% Rule: keep only relevant work. Drop unrelated jobs (like retail) once you
%% have projects and team or lab work to fill the page.
\section{Experience}

\entry{Power Subsystem Engineer, State University CubeSat Team}{Fall 24 - Spring 26}
\begin{itemize}
  \item Cut power board mass by 30 g by replacing 6 linear regulators with one buck converter board, which then passed thermal vacuum testing from -20 \textdegree{}C to 60 \textdegree{}C on the first attempt with no board rework needed
  \item Found 2 battery protection faults before launch review by writing a Python test script that runs 200 charge and discharge cycles overnight and flags any cell that drifts more than 20 mV from the pack average
\end{itemize}
\vspace{\itemgap}

\entry{Undergraduate Research Assistant, Power Electronics Lab, State University}{Summer 25}
\begin{itemize}
  \item Reduced GaN switching loss measurement error from 12\% to 3\% by adding probe deskew calibration to the double-pulse test setup, which the lab now uses as the standard setup for every GaN device it tests
  \item Cut test time per device from 40 to 8 minutes by automating oscilloscope capture with Python and PyVISA
\end{itemize}

%% ---------------- Education ----------------
%% Rules: degree on top, school under it, GPA on the degree line.
%% Bold the "Relevant coursework:" and "Activities:" labels so eyes find them.
%% No period after the last item.
\section{Education}

\entry{B.S. Electrical Engineering, GPA: 3.62/4.00}{Expected Spring 27}
State University\par
\textbf{Relevant coursework:} Signals and Systems, Electronics I and II, Power Electronics, Digital Design, Embedded Systems, Control Systems, Electromagnetics, Probability for Engineers, Microcontroller Lab\par
\textbf{Activities:} Tau Beta Pi, IEEE student branch\par

%% ---------------- Skills ----------------
%% Rules: one labeled line per category, most important first, only what you
%% can prove. Never two columns. No soft skills ("communication", "problem
%% solving"). Knowledge areas like orbital mechanics belong in coursework
%% unless they are certified. Add a "Hardware:" line (scope, logic analyzer,
%% rework station) if you have bench tools to list.
\section{Skills}

\textbf{Certifications:} CSWA (Certified SolidWorks Associate), IPC J-STD-001 soldering\par
\textbf{Software:} KiCad, LTspice, MATLAB, Vivado, Python, C, Verilog, SolidWorks, Git\par

\end{document}
